\section{Overview}\label{overview}

In this first chapter we will focus on the physical appearance of the
LXR, describing the front panel controls as well as the connection jacks
on the back. The basic menu navigation is also explained.

\subsection{The Front}\label{the-front}

\begin{figure}[H]
  \centering
  \includegraphics[width=\linewidth]{front_panel.png}
  \caption{LXR front panel with numbered control regions.}
\end{figure}

\subsubsection{1) Display}\label{display}

The display shows parameter values from the selected menu page.

\subsubsection{2) Knobs}\label{knobs}

The 4 knobs edit the values shown in the display. A parameter fetch
option --- which prevents value jumps when switching menu pages --- can
be activated in the settings menu.

\subsubsection{3) Record/Play Buttons}\label{recordplay-buttons}

\textbf{The Play button} starts and stops playback. The current pattern
is reset to the first step when the sequencer is stopped. {[}Alternative
SHIFT function: reset last sound parameters loaded from file.{]}

\textbf{The Record button} activates the recording function. When the
LED is lit, all incoming MIDI notes are recorded into the active
pattern. Knob movements are also recorded to the selected automation
track. {[}Alternative SHIFT function: record settings.{]}

\subsubsection{4) Shift Button}\label{shift-button}

Activates an alternate function set for the buttons. Hold SHIFT while
pressing another button to activate the alternative function. A global
setting can change the SHIFT button to work as a toggle instead, in
which case the led indicates whether SHIFT is active or not.

\subsubsection{4) Copy/Clear Button}\label{copyclear-button}

Used to copy and {[}clear (alternative SHIFT function){]} sequencer
data.

\subsubsection{5) 16 Sequencer Buttons}\label{sequencer-buttons}

Function depends on the active mode:

\begin{itemize}
\tightlist
\item
  \textbf{Voice mode:} set/reset the 16 main steps in the sequencer.
\item
  \textbf{Performance mode:} buttons 1--7 trigger manual rolls. Buttons
  9--16 control the looper.
\item
  \textbf{Step Edit mode:} select one of the 16 main steps for editing.
\item
  \textbf{Menu mode:} no function.
\end{itemize}

\subsubsection{6) Encoder}\label{encoder}

Used to navigate through the menu. Press the encoder to activate a
selected option or enter a single-parameter view. {[}Hold shift while in
single-parameter view to view and edit the \textbf{morph} parameter
where applicable.{]}

\subsubsection{7) 8 Select Buttons}\label{select-buttons}

Function depends on the active mode:

\begin{itemize}
\tightlist
\item
  \textbf{Voice Edit mode:} select the synth section to edit
  (Oscillators, Mix, Envelopes, etc.)
\item
  \textbf{Performance mode:} change the playing pattern.
\item
  \textbf{Step Edit mode:} select/edit sub steps.
\item
  \textbf{Menu/Pattern generator mode:} no function.
\end{itemize}

\subsubsection{8) 4 Mode Buttons}\label{mode-buttons}

Switch between the 4 operating modes. Modes shown in brackets are
selected using the shift button.

\begin{itemize}
\tightlist
\item
  Button 1 --- Voice edit mode
\item
  Button 2 --- Performance mode {[}Pattern generator mode{]}
\item
  Button 3 --- Step Edit mode
\item
  Button 4 --- Load/Save {[}Settings menu{]}
\end{itemize}

\subsubsection{9) 7 Voice Buttons}\label{voice-buttons}

The voice buttons select the active track to edit. Together with the
shift key they mute/unmute the voices. They behave the same in all
operation modes except performance mode --- see the performance mode
chapter for details.

\subsection{The Back}\label{the-back}

\begin{figure}[H]
  \centering
  \includegraphics[width=0.85\linewidth]{back_panel.jpeg}
  \caption{LXR rear panel connections.}
\end{figure}

\subsubsection{USB}\label{usb}

The synthesizer provides a class-compliant USB MIDI interface.

\subsubsection{Memory}\label{memory}

Insert an SD card to load and save preset and pattern data. The card
must be inserted with the label facing downward.

\attention{The SD card must use a FAT32 filesystem for reliable operation. The synth will work
with FAT16 cards too, but \textbf{firmware updates are only possible from FAT32 cards}.}

\subsubsection{4 Audio Outputs}\label{audio-outputs}

4 mono audio outputs at line level. Can also be used as 2 stereo output
pairs.

\subsubsection{MIDI}\label{midi}

MIDI in and output via standard DIN connectors.

\subsubsection{Power Connector}\label{power-connector}

\textbf{7--9V DC} power plug, \textbf{centre pin positive, at least 600
mA}.

\subsection{Menu Navigation}\label{menu-navigation}

The upper row of the display shows the short name of the parameter. The
bottom row shows its current value. Use one of the 4 knobs to change it.

\oled{frq \quad Res \quad typ \quad drv >}{32 \quad\; 120 \quad LP \quad\;\; 0}

\begin{itemize}
\tightlist
\item
  The encoder selects a menu parameter. A capital letter at the
  beginning of the name indicates the selected parameter. In the example
  above, \texttt{Res} is selected.
\item
  If there are more than 4 entries in the current menu, a
  \texttt{\textgreater{}} or \texttt{\textless{}} sign in the
  upper-right corner indicates a second page. The two pages toggle by
  pressing the menu button again.
\end{itemize}

Pushing the encoder opens the detail page:

\oled{\textbf{Filter}}{Resonance \hfill 120}

On the detail page you can see the full parameter name and change its
value with the encoder. This is useful for fine adjustments where the
knob control is too coarse. When viewing a drumkit parameter in VOICE
mode, hold SHIFT in this mode to view and edit the corresponding
\textbf{morph} parameter. Push the encoder again to return to the normal
menu mode.
